% TOP_PROBLEM: {"id": 7900015, "title": "One-Dimensional Lieb-Thirring Constants", "category_id": 16, "category": {"id": 16, "name": "physics", "display_name": "Mathematical Physics", "description": "Problems at the intersection of mathematics and physics.", "slug": "physics", "order_index": 16, "created_at": "2026-07-31T15:26:25.671Z"}, "status": "partially_solved"}
% FIRST_POSTED: 2026-10-06
% ATTEMPT_STATUS: solved
% OPENAI_CLAIM: {"schema_version":1,"source_repo":"https://github.com/openai/math","source_commit":"adc7f1241b42e322a6451854ab7e4b4c146bf78a","release_date":"2026-10-06","match":"stronger","resolution":"proved","independently_reviewed":false,"second_pass":"agreed","families":[{"family":"262","title":"Sharp finite-matrix Lieb–Thirring inequalities and all equality cases","subject":"Mathematical physics","manuscripts":[{"title":"Equality cases in the sharp one-dimensional matrix Lieb–Thirring inequality","pdf_path":"preprints/Equality-cases-in-the-sharp-one-dimensional-matrix-Lieb-Thirring-inequality-October-5-2026/sharp-one-dimensional-lieb-thirring-inequalities-matrix-potentials.pdf","readme_path":"preprints/Equality-cases-in-the-sharp-one-dimensional-matrix-Lieb-Thirring-inequality-October-5-2026/README.md","date":"2026-10-05","theorem_ref":"thm:main (sections/introduction.tex:34)","statement_sources":[{"path":"preprints/Equality-cases-in-the-sharp-one-dimensional-matrix-Lieb-Thirring-inequality-October-5-2026/build/sections/introduction.tex","line":34,"label":"thm:main"}]},{"title":"Sharp one-dimensional Lieb–Thirring constants","pdf_path":"preprints/Sharp-One-Dimensional-Lieb-Thirring-Constants-September-23-2026/paper.pdf","readme_path":"preprints/Sharp-One-Dimensional-Lieb-Thirring-Constants-September-23-2026/README.md","date":"2026-09-23","theorem_ref":"thm:main (sections/introduction.tex:37)","statement_sources":[{"path":"preprints/Sharp-One-Dimensional-Lieb-Thirring-Constants-September-23-2026/build/sections/introduction.tex","line":37,"label":"thm:main"}]},{"title":"Sharp one-dimensional Lieb–Thirring inequalities for matrix potentials","pdf_path":"preprints/Sharp-one-dimensional-Lieb-Thirring-inequalities-for-matrix-potentials-October-5-2026/sharp-matrix-lieb-thirring.pdf","readme_path":"preprints/Sharp-one-dimensional-Lieb-Thirring-inequalities-for-matrix-potentials-October-5-2026/README.md","date":"2026-10-05","theorem_ref":"thm:main (sections/introduction.tex:28)","statement_sources":[{"path":"preprints/Sharp-one-dimensional-Lieb-Thirring-inequalities-for-matrix-potentials-October-5-2026/build/sections/introduction.tex","line":28,"label":"thm:main"}]}],"lean":{"doc_path":"lean/docs/262.md","comparators":["lean/ComparatorChallenges/LiebThirring.json","lean/ComparatorChallenges/LiebThirring.lean"],"declarations":["OAI.SharpLiebThirring.sharp_lieb_thirring"],"covers_main_theorem":false},"reasoning_trace":null}]}
\documentclass[11pt]{article}
\usepackage[margin=25mm]{geometry}
\usepackage{fontspec}
\setmainfont{Latin Modern Roman}
\usepackage{amsmath,amssymb,mathtools}
\usepackage{unicode-math}
\setmathfont{Latin Modern Math}
\usepackage{enumitem}
\usepackage{xurl,hyperref,bookmark}
\hypersetup{colorlinks=true,urlcolor=blue}
\setcounter{secnumdepth}{0}
\setlength{\parindent}{0pt}
\setlength{\parskip}{5pt}
\setlength{\emergencystretch}{3em}
\begin{document}
\section*{One-Dimensional Lieb-Thirring Constants}

\subsection{OpenAI's claim}
The linked manuscripts present the following claims by OpenAI: All equality cases of that matrix inequality are claimed to be channels in a constant unitary basis, each either zero or a scalar single-state soliton; The scalar one-dimensional Lieb--Thirring optimal constant is claimed to be the one-bound-state value for every $1/2<\gamma<3/2$; The same sharp constant is claimed for Hermitian positive matrix potentials of every finite size, even when values at different points fail to commute. Together they give a claimed affirmative answer in the matching scope.

\subsection{Mathematical statement of the claimed result}
\paragraph{Result 1.} Fix a finite matrix size \(m\ge1\), \(1/2<\gamma<3/2\), and \(p=\gamma+1/2\). For measurable Hermitian \(W(x)\ge0\) with \(\int_{\mathbb R}\operatorname{tr}(W^p)<\infty\), form \(H_W=-d^2/dx^2-W\) by its \(H^1\) quadratic form. The source claims
\[
 \sum_{\lambda_j(H_W)<0}|\lambda_j|^\gamma
 \le C_\gamma\int_{\mathbb R}\operatorname{tr}(W^p),\quad
 C_\gamma=\left(\frac{\gamma-1/2}{\gamma+1/2}\right)^{\gamma-1/2}
 \frac{\Gamma(\gamma+1)}{\sqrt\pi\Gamma(\gamma+3/2)}.
\]
Writing \(r=(\gamma-1/2)^{-1}\), equality occurs exactly for a constant unitary \(U\), \(0\le k\le m\), \(a_j>0\), \(x_j\in\mathbb R\), and
\[
 W(x)=U\operatorname{diag}(w_1(x),\ldots,w_k(x),0,\ldots,0)U^*,\quad
 w_j(x)=(r+1)a_j^2\operatorname{sech}^2(ra_j(x-x_j))
\]
almost everywhere. Each nonzero channel has just the negative eigenvalue \(-a_j^2\), and \(k=0\) means the zero potential.

\paragraph{Result 2.} For \(1/2<\gamma<3/2\) and \(W\ge0\) in \(L^{\gamma+1/2}(\mathbb R)\), put \(H_{-W}=-d^2/dx^2-W\). Counting negative eigenvalues with multiplicity, the source claims
\[
 \sum_{\lambda_j<0}|\lambda_j|^\gamma\le C_\gamma\int_{\mathbb R}W^{\gamma+1/2},\quad
 C_\gamma=\left(\frac{\gamma-1/2}{\gamma+1/2}\right)^{\gamma-1/2}
 \frac{\Gamma(\gamma+1)}{\sqrt\pi\Gamma(\gamma+3/2)}.
\]
This constant is optimal. An equality potential is \(W(x)=(r+1)\operatorname{sech}^2(rx)\), \(r=(\gamma-1/2)^{-1}\).

\paragraph{Result 3.} Let \(m\ge1\) be finite, \(1/2<\gamma<3/2\), and \(W:\mathbb R\to\mathbb C^{m\times m}\) be measurable, Hermitian, and positive semidefinite, with \(\int\operatorname{tr}(W^{\gamma+1/2})<\infty\). For the form-defined operator \(H_W=-d^2/dx^2-W\), the claim is
\[
 \sum_{\lambda_j(H_W)<0}|\lambda_j|^\gamma\le C_\gamma\int_{\mathbb R}\operatorname{tr}(W^{\gamma+1/2}),\quad
 C_\gamma=\left(\frac{\gamma-1/2}{\gamma+1/2}\right)^{\gamma-1/2}
 \frac{\Gamma(\gamma+1)}{\sqrt\pi\Gamma(\gamma+3/2)}.
\]
The constant is sharp for each \(m\), witnessed by \(W(x)=\operatorname{diag}((r+1)\operatorname{sech}^2(rx),0,\ldots,0)\), \(r=(\gamma-1/2)^{-1}\).

\subsection{Outline of the argument}
The action estimate is refined to keep each convexity and trace defect. Vanishing defects force pointwise matrix relations and fixed invariant channels, after which the scalar differential equations identify the independent soliton scales and centers.

A finite-interval action inequality is obtained through a regularized matrix field and a rank-one comparison. Matrix paths with prescribed endpoints bound the action, which is integrated into the full negative-eigenvalue moment inequality; scalar solitons provide sharpness.

An integrated trace inequality controls a regularized matrix path without simultaneous diagonalization. The endpoint action estimate gives the complete negative-spectrum moment bound, and embedding a scalar one-bound-state channel shows the dimension-independent constant cannot improve.

\subsection{Scope relative to this problem}
The equality-classification theorem proves the optimal one-bound-state constant for every 1/2<gamma<3/2, includes the scalar m=1 target, and further characterizes all finite-matrix extremizers. Its integrability class is the full L\textasciicircum{}(gamma+1/2) nonnegative potential class. For every 1/2<gamma<3/2 and every nonnegative L\textasciicircum{}(gamma+1/2) scalar potential, the theorem identifies the full negative-eigenvalue moment constant with the predicted one-bound-state value and exhibits an extremizer. No truncated spectral sum or smaller exponent range is used. The matrix theorem proves the optimal predicted one-dimensional constant for every 1/2<gamma<3/2 and every finite matrix dimension m; scalar m=1 is exactly the problem statement target, while matrix dimensions are additional scope. The claimed conclusion supplies the requested case and extends it in the directions just specified.

\subsection{Formal verification}
For family 262, the documented result is the sharp scalar one-dimensional inequality in the intermediate exponent range, including all negative eigenvalues. The matrix and full equality-classification companions are not selected. The listed Comparator files are \texttt{LiebThirring.lean}. Named declarations include \texttt{OAI.SharpLiebThirring.sharp\_lieb\_thirring}. This site has not executed the Lean code. The scope descriptions do not certify the remaining manuscript arguments.

\subsection{Status caveat}
OpenAI attributes these manuscripts to work by an internal model. Its release cautions that arguments without formalization can contain mistakes. This catalogue has not independently checked the proofs; the labels report the mathematical claims attributed to OpenAI.

\subsection{Sources}
\begin{itemize}
\item OpenAI, \emph{Equality cases in the sharp one-dimensional matrix Lieb–Thirring inequality}, 2026-10-05; citation key \texttt{OAI:Equality-cases-in-the-sharp-one-dimensional-matrix-Lieb-Thirring-inequality-October-5-2026}. \href{https://github.com/openai/math/blob/adc7f1241b42e322a6451854ab7e4b4c146bf78a/preprints/Equality-cases-in-the-sharp-one-dimensional-matrix-Lieb-Thirring-inequality-October-5-2026/sharp-one-dimensional-lieb-thirring-inequalities-matrix-potentials.pdf}{PDF (GitHub)}; \href{https://github.com/openai/math/raw/adc7f1241b42e322a6451854ab7e4b4c146bf78a/preprints/Equality-cases-in-the-sharp-one-dimensional-matrix-Lieb-Thirring-inequality-October-5-2026/sharp-one-dimensional-lieb-thirring-inequalities-matrix-potentials.pdf}{PDF (direct)}; \href{https://github.com/openai/math/blob/adc7f1241b42e322a6451854ab7e4b4c146bf78a/preprints/Equality-cases-in-the-sharp-one-dimensional-matrix-Lieb-Thirring-inequality-October-5-2026/README.md}{Manuscript page and citation}; \href{https://github.com/openai/math/blob/main/preprints/Equality-cases-in-the-sharp-one-dimensional-matrix-Lieb-Thirring-inequality-October-5-2026/sharp-one-dimensional-lieb-thirring-inequalities-matrix-potentials.pdf}{Latest version}; Statement source: \href{https://github.com/openai/math/blob/adc7f1241b42e322a6451854ab7e4b4c146bf78a/preprints/Equality-cases-in-the-sharp-one-dimensional-matrix-Lieb-Thirring-inequality-October-5-2026/build/sections/introduction.tex\#L34}{thm:main}.
\item OpenAI, \emph{Sharp one-dimensional Lieb–Thirring constants}, 2026-09-23; citation key \texttt{OAI:Sharp-One-Dimensional-Lieb-Thirring-Constants-September-23-2026}. \href{https://github.com/openai/math/blob/adc7f1241b42e322a6451854ab7e4b4c146bf78a/preprints/Sharp-One-Dimensional-Lieb-Thirring-Constants-September-23-2026/paper.pdf}{PDF (GitHub)}; \href{https://github.com/openai/math/raw/adc7f1241b42e322a6451854ab7e4b4c146bf78a/preprints/Sharp-One-Dimensional-Lieb-Thirring-Constants-September-23-2026/paper.pdf}{PDF (direct)}; \href{https://github.com/openai/math/blob/adc7f1241b42e322a6451854ab7e4b4c146bf78a/preprints/Sharp-One-Dimensional-Lieb-Thirring-Constants-September-23-2026/README.md}{Manuscript page and citation}; \href{https://github.com/openai/math/blob/main/preprints/Sharp-One-Dimensional-Lieb-Thirring-Constants-September-23-2026/paper.pdf}{Latest version}; Statement source: \href{https://github.com/openai/math/blob/adc7f1241b42e322a6451854ab7e4b4c146bf78a/preprints/Sharp-One-Dimensional-Lieb-Thirring-Constants-September-23-2026/build/sections/introduction.tex\#L37}{thm:main}.
\item OpenAI, \emph{Sharp one-dimensional Lieb–Thirring inequalities for matrix potentials}, 2026-10-05; citation key \texttt{OAI:Sharp-one-dimensional-Lieb-Thirring-inequalities-for-matrix-potentials-October-5-2026}. \href{https://github.com/openai/math/blob/adc7f1241b42e322a6451854ab7e4b4c146bf78a/preprints/Sharp-one-dimensional-Lieb-Thirring-inequalities-for-matrix-potentials-October-5-2026/sharp-matrix-lieb-thirring.pdf}{PDF (GitHub)}; \href{https://github.com/openai/math/raw/adc7f1241b42e322a6451854ab7e4b4c146bf78a/preprints/Sharp-one-dimensional-Lieb-Thirring-inequalities-for-matrix-potentials-October-5-2026/sharp-matrix-lieb-thirring.pdf}{PDF (direct)}; \href{https://github.com/openai/math/blob/adc7f1241b42e322a6451854ab7e4b4c146bf78a/preprints/Sharp-one-dimensional-Lieb-Thirring-inequalities-for-matrix-potentials-October-5-2026/README.md}{Manuscript page and citation}; \href{https://github.com/openai/math/blob/main/preprints/Sharp-one-dimensional-Lieb-Thirring-inequalities-for-matrix-potentials-October-5-2026/sharp-matrix-lieb-thirring.pdf}{Latest version}; Statement source: \href{https://github.com/openai/math/blob/adc7f1241b42e322a6451854ab7e4b4c146bf78a/preprints/Sharp-one-dimensional-Lieb-Thirring-inequalities-for-matrix-potentials-October-5-2026/build/sections/introduction.tex\#L28}{thm:main}.
\item Family 262: \href{https://github.com/openai/math/blob/adc7f1241b42e322a6451854ab7e4b4c146bf78a/lean/docs/262.md}{Lean scope}.
\item \href{https://github.com/openai/math/blob/adc7f1241b42e322a6451854ab7e4b4c146bf78a/lean/ComparatorChallenges/LiebThirring.json}{Comparator LiebThirring.json}.
\item \href{https://github.com/openai/math/blob/adc7f1241b42e322a6451854ab7e4b4c146bf78a/lean/ComparatorChallenges/LiebThirring.lean}{Comparator LiebThirring.lean}.
\item \href{https://github.com/openai/math/blob/adc7f1241b42e322a6451854ab7e4b4c146bf78a/CONTENTS.md}{OpenAI catalogue entry}.
\end{itemize}
\end{document}
